%----------------------------------------------------------------------------------------
% RADIALEN EN DE EENHEIDSCIRKEL. Alle opgaven zijn eigen werk.
%
% Opbouw: Start, leerstof (theorie, werkwijze of uitleg, voorbeeld, korte oefening), en op een
% nieuwe pagina de Opdrachten (door elkaar, stijgende moeilijkheid), Koppel en Test jezelf.
% Labels van dit hoofdstuk beginnen met h1-.
%----------------------------------------------------------------------------------------

\mychapter{1}{De eenheidscirkel}

\lkr{Leerplan B+S: LPD 13-16 (goniometrische functies), het deel over radialen, georiënteerde hoeken en de goniometrische getallen op de goniometrische cirkel. Uit de 2de graad kennen leerlingen sinus, cosinus en tangens in een rechthoekige driehoek, in graden. Laat de bijzondere waarden niet van buiten leren zonder figuur: de twee tekendriehoeken en de eenheidscirkel volstaan om ze altijd terug te vinden.}

In de 2de graad waren sinus, cosinus en tangens verhoudingen van zijden in een rechthoekige driehoek, voor hoeken tussen $0^\circ$ en $90^\circ$. Een rad, een golf of een slinger draait echter verder dan $90^\circ$, en soms in tegenwijzerzin. Daarom breiden we de goniometrische getallen uit tot \textit{alle} hoeken, met de eenheidscirkel. En we meten hoeken voortaan in radialen.

\begin{oefeningen}[Start]
\oef{h1-start-1}

\oef{h1-verken}
\end{oefeningen}

%========================================================================================
\section{Radialen}

\naast{%
\begin{theorie}
Een hoek van \textbf{\textcolor{accent}{1 radiaal}} is de middelpuntshoek die op een cirkel met straal $r$ een boog met lengte $r$ afsnijdt. Een volle cirkel is dus $2\pi$ radialen:
\[\formule{180^\circ = \pi \text{ rad}}.\]
Een middelpuntshoek van $\theta$ rad op een cirkel met straal $r$ snijdt af:
\[\formule{\text{booglengte } s = r\theta} \qquad \formule{\text{sectoroppervlakte } A = \tfrac12 r^2\theta}.\]
\end{theorie}}{%
\begin{tikzpicture}[scale=1.25,font=\scriptsize]
  \fill[accent!15] (0,0) -- (0:1) arc (0:57.3:1) -- cycle;
  \draw[grijs] (0,0) circle (1);
  \draw[accent,thick] (0:1) arc (0:57.3:1);
  \draw (0,0) -- node[below] {$r$} (1,0);
  \draw (0,0) -- node[above left] {$r$} (57.3:1);
  \node[accent] at (28:1.2) {$r$};
  \node at (25:.4) {$1$ rad};
\end{tikzpicture}}

Radialen zijn de natuurlijke maat voor hoeken: met radialen worden de formules voor booglengte en afgeleiden eenvoudig. Vanaf nu meten we hoeken in radialen, tenzij er een gradenteken staat. Een hoek van $1$ rad is ongeveer $57{,}3^\circ$.

\voorbeeld{}{Zet $135^\circ$ om naar radialen, en $2{,}5$ rad naar graden.\label{vb:h1-radialen}}{%
\uitleg{Vermenigvuldig met $\frac{\pi}{180}$ om naar radialen te gaan, met $\frac{180}{\pi}$ om terug te gaan.}
\begin{stappenbreed}
135^\circ = 135 \cdot \frac{\pi}{180} = \frac{3\pi}{4} & Exact, met $\pi$. \\
2{,}5 \text{ rad} = 2{,}5 \cdot \frac{180^\circ}{\pi} \approx 143{,}2^\circ & Benadering.
\end{stappenbreed}}

\voorbeeld{}{Een pizzapunt heeft straal $15$ cm en een tophoek van $0{,}8$ rad. Hoe lang is de korst, en hoe groot is de punt?\label{vb:h1-boog}}{%
\uitleg{De korst is de boog, de punt is een cirkelsector.}
\begin{stappenbreed}
s = r\theta = 15 \cdot 0{,}8 = 12 & Booglengte in cm. \\
A = \tfrac12 r^2 \theta = \tfrac12 \cdot 225 \cdot 0{,}8 = 90 & Oppervlakte in cm$^2$.
\end{stappenbreed}}

\begin{oefeningen}
\oef{h1-radialen}

\oef{h1-sector}
\end{oefeningen}

%========================================================================================
\section{Sinus en cosinus op de eenheidscirkel}

De \textbf{eenheidscirkel} heeft middelpunt $O(0, 0)$ en straal $1$. Een \textbf{georiënteerde hoek} $\theta$ meten we vanaf de positieve $x$-as: in tegenwijzerzin positief, in wijzerzin negatief. Zo hoort bij elk reëel getal $\theta$ één punt $P$ op de cirkel. Hoeken die een veelvoud van $2\pi$ verschillen, horen bij hetzelfde punt.

\naast{%
\begin{theorie}
Is $P$ het punt van de eenheidscirkel bij de hoek $\theta$, dan is
\[\formule{\cos\theta = x_P} \qquad \formule{\sin\theta = y_P} \qquad \formule{\tan\theta = \frac{\sin\theta}{\cos\theta}} \ (\cos\theta \neq 0).\]
\begin{stelling}{Grondformule}
Voor elke $\theta$: \quad $\formule{\sin^2\theta + \cos^2\theta = 1}$.
\end{stelling}
\end{theorie}}{%
\begin{tikzpicture}[scale=1.5,font=\scriptsize]
  \eenheidscirkel
  \draw[accent,thick] (0,0) -- (130:1);
  \draw[accent,dashed] (130:1) -- ({cos(130)},0);
  \draw[accent,dashed] (130:1) -- (0,{sin(130)});
  \draw[->,accent] (.25,0) arc (0:130:.25);
  \node[accent] at (60:.38) {$\theta$};
  \fill[accent] (130:1) circle (.03) node[above left] {$P$};
  \node[below] at ({cos(130)},0) {$\cos\theta$};
  \node[right] at (0,{sin(130)}) {$\sin\theta$};
  \node[below right] at (1,0) {$1$};
\end{tikzpicture}}

Voor een hoek tussen $0$ en $\frac{\pi}{2}$ zijn dit de oude definities: in de rechthoekige driehoek met schuine zijde $1$ is de overstaande zijde $\sin\theta$ en de aanliggende $\cos\theta$. De grondformule is de stelling van Pythagoras in die driehoek. Het teken van $\cos\theta$ en $\sin\theta$ hangt af van het kwadrant waarin $P$ ligt.

\begin{theorie}
\begin{stelling}{Bijzondere waarden}
\renewcommand{\arraystretch}{1.6}
\begin{tabular}{@{}c|*{5}{>{\centering\arraybackslash}p{1.25cm}}@{}}
$\theta$ & $0$ & $\tfrac{\pi}{6}$ & $\tfrac{\pi}{4}$ & $\tfrac{\pi}{3}$ & $\tfrac{\pi}{2}$ \\ \hline
$\sin\theta$ & $\formule{0}$ & $\formule{\tfrac12}$ & $\formule{\tfrac{\sqrt2}{2}}$ & $\formule{\tfrac{\sqrt3}{2}}$ & $\formule{1}$ \\
$\cos\theta$ & $\formule{1}$ & $\formule{\tfrac{\sqrt3}{2}}$ & $\formule{\tfrac{\sqrt2}{2}}$ & $\formule{\tfrac12}$ & $\formule{0}$ \\
$\tan\theta$ & $\formule{0}$ & $\formule{\tfrac{\sqrt3}{3}}$ & $\formule{1}$ & $\formule{\sqrt3}$ & bestaat niet
\end{tabular}
\end{stelling}
\end{theorie}

\voorbeeld{}{Bepaal exact $\sin\frac{5\pi}{6}$, $\cos\frac{5\pi}{6}$ en $\tan\frac{5\pi}{6}$.\label{vb:h1-waarden}}{%
\naast{%
\uitleg{$\frac{5\pi}{6} = \pi - \frac{\pi}{6}$: het punt is het spiegelbeeld van het punt bij $\frac{\pi}{6}$ om de $y$-as. Dezelfde hoogte, tegengestelde $x$.}
\begin{stappenvrij}{4cm}{3.6cm}
\sin\tfrac{5\pi}{6} = \sin\tfrac{\pi}{6} = \tfrac12 & Zelfde $y$. \\
\cos\tfrac{5\pi}{6} = -\cos\tfrac{\pi}{6} = -\tfrac{\sqrt3}{2} & Tegengestelde $x$. \\
\tan\tfrac{5\pi}{6} = -\tfrac{1}{\sqrt3} = -\tfrac{\sqrt3}{3} & $\sin$ gedeeld door $\cos$.
\end{stappenvrij}}{%
\begin{tikzpicture}[scale=1.4,font=\scriptsize]
  \eenheidscirkel
  \draw[accent] (0,0) -- (30:1) (0,0) -- (150:1);
  \draw[accent,dashed] (30:1) -- (150:1);
  \fill[accent] (30:1) circle (.03) node[right] {$\frac{\pi}{6}$};
  \fill[accent] (150:1) circle (.03) node[left] {$\frac{5\pi}{6}$};
\end{tikzpicture}}}

\begin{oefeningen}
\oef{h1-punt}
\end{oefeningen}

\voorbeeld{}{Gegeven: $\sin\theta = \frac35$ en $\frac{\pi}{2} < \theta < \pi$. Bepaal $\cos\theta$ en $\tan\theta$.\label{vb:h1-grondformule}}{%
\uitleg{De grondformule geeft $\cos\theta$ op het teken na. In het tweede kwadrant is de $x$-coördinaat negatief.}
\begin{stappenbreed}
\cos^2\theta = 1 - \tfrac{9}{25} = \tfrac{16}{25} & Grondformule. \\
\cos\theta = -\tfrac45 & Tweede kwadrant: negatief. \\
\tan\theta = \dfrac{3/5}{-4/5} = -\tfrac34 &
\end{stappenbreed}}

\begin{oefeningen}
\oef{h1-waarden}

\oef{h1-grondf}
\end{oefeningen}

%========================================================================================
\section{Verwante hoeken}

Spiegel je het punt $P$ bij $\theta$ om een as of om de oorsprong, dan krijg je punten bij hoeken die met $\theta$ verwant zijn. Hun coördinaten zijn dezelfde, op het teken of de volgorde na.

\begin{theorie}
\begin{stelling}{Verwante hoeken}
\renewcommand{\arraystretch}{1.3}
\begin{tabular}{@{}lll@{}}
tegengestelde hoeken & $\formule{\sin(-\theta) = -\sin\theta}$ & $\formule{\cos(-\theta) = \cos\theta}$ \\
supplementaire hoeken & $\formule{\sin(\pi - \theta) = \sin\theta}$ & $\formule{\cos(\pi - \theta) = -\cos\theta}$ \\
antisupplementaire hoeken & $\formule{\sin(\pi + \theta) = -\sin\theta}$ & $\formule{\cos(\pi + \theta) = -\cos\theta}$ \\
complementaire hoeken & $\formule{\sin\bigl(\tfrac{\pi}{2} - \theta\bigr) = \cos\theta}$ & $\formule{\cos\bigl(\tfrac{\pi}{2} - \theta\bigr) = \sin\theta}$
\end{tabular}
\end{stelling}
\end{theorie}

\figuren{%
\begin{tikzpicture}[scale=1.25,font=\scriptsize]
  \eenheidscirkel
  \foreach \h/\n/\p in {35/\theta/right,145/\pi - \theta/left,215/\pi + \theta/left,-35/-\theta/right}{
    \draw[accent] (0,0) -- (\h:1);
    \fill[accent] (\h:1) circle (.03) node[\p] {$\n$};}
  \draw[accent,dashed] (35:1) -- (145:1) (35:1) -- (-35:1) (35:1) -- (215:1);
\end{tikzpicture}\hspace{1cm}
\begin{tikzpicture}[scale=1.25,font=\scriptsize]
  \eenheidscirkel
  \draw[grijs,dashed] (-1.1,-1.1) -- (1.2,1.2);
  \foreach \h/\n/\p in {25/\theta/right,65/\frac{\pi}{2} - \theta/above right}{
    \draw[accent] (0,0) -- (\h:1);
    \fill[accent] (\h:1) circle (.03) node[\p] {$\n$};}
\end{tikzpicture}}

Supplementair: spiegelen om de $y$-as. Tegengesteld: om de $x$-as. Antisupplementair: om de oorsprong. Complementair: om de rechte $y = x$, zodat $x$ en $y$ wisselen. De tangens volgt telkens als quotiënt; zo is $\tan(\pi + \theta) = \tan\theta$.

\voorbeeld{}{Bepaal exact $\cos\frac{4\pi}{3}$.\label{vb:h1-verwant}}{%
\uitleg{$\frac{4\pi}{3} = \pi + \frac{\pi}{3}$: een antisupplementaire hoek van $\frac{\pi}{3}$.}
\begin{stappenbreed}
\cos\tfrac{4\pi}{3} = \cos\bigl(\pi + \tfrac{\pi}{3}\bigr) = -\cos\tfrac{\pi}{3} & Antisupplementair. \\
= -\tfrac12 & Bijzondere waarde.
\end{stappenbreed}}

\begin{werkwijze}
Trek veelvouden van $2\pi$ af tot de hoek tussen $-\pi$ en $\pi$ (of tussen $0$ en $2\pi$) ligt. Schrijf hem als $\pi - \alpha$, $\pi + \alpha$ of $-\alpha$ met $\alpha$ tussen $0$ en $\frac{\pi}{2}$. Gebruik de tabel van de verwante hoeken, en controleer het teken met het kwadrant.
\end{werkwijze}

\nujij{Bepaal exact $\sin\bigl(-\frac{7\pi}{4}\bigr)$.}{$-\frac{7\pi}{4} + 2\pi = \frac{\pi}{4}$, dus $\sin\bigl(-\frac{7\pi}{4}\bigr) = \sin\frac{\pi}{4} = \frac{\sqrt2}{2}$.}

\begin{oefeningen}
\oef{h1-verwant}

\oef{h1-herleid}
\end{oefeningen}

%========================================================================================
\opdrachtenkop

\begin{oefeningen}
\needspace{10\baselineskip}
\oef{h1-o1}

\needspace{12\baselineskip}
\oef{h1-o10}

\oef{h1-o13}

\needspace{10\baselineskip}
\oef{h1-o2}

\needspace{12\baselineskip}
\oef{h1-o11}

\needspace{12\baselineskip}
\oef{h1-o14}

\oef{h1-o3}

\oef{h1-o15}

\oef{h1-o12}

\oef{h1-o4}

\oef{h1-o16}

\oef{h1-o5}

\needspace{12\baselineskip}
\oef{h1-o17}

\oef{h1-o7}

\needspace{10\baselineskip}
\oef{h1-o20}

\oef{h1-o6}

\oef{h1-o22}

\oef{h1-o19}

\needspace{14\baselineskip}
\oef{h1-o18}

\oef{h1-o23}

\needspace{12\baselineskip}
\oef{h1-o21}

\needspace{14\baselineskip}
\oef{h1-o24}

\oef{h1-o8}

\oef{h1-o26}

\needspace{16\baselineskip}
\oef{h1-o25}

\oef{h1-o9}
\end{oefeningen}

%----------------------------------------------------------------------------------------
% KOPPEL: verbanden met andere leerstof
%----------------------------------------------------------------------------------------
\needspace{20\baselineskip}
\opdrachtsectie{Koppel}

\lkr{Algemene functies en Matrices: kijk na wat al gezien is. Het pakket legt de volgorde niet vast.}

\begin{oefeningen}

\oef{h1-koppel-2}

\oef{h1-koppel-3}
\end{oefeningen}

%----------------------------------------------------------------------------------------
\opdrachtsectie{Test jezelf}

\begin{oefeningen}
\oef{h1-t1}

\oef{h1-t2}

\oef{h1-t3}

\oef{h1-t4}

\oef{h1-t5}

\oef{h1-t6}

\oef{h1-t7}

\oef{h1-t8}
\end{oefeningen}

\herhaalschema{%
  \fv{oef:h1-t1}{vb:h1-radialen}{oef:h1-radialen}%
  \fv{oef:h1-t2}{vb:h1-boog}{oef:h1-sector}%
  \fv{oef:h1-t3}{vb:h1-waarden}{oef:h1-waarden}%
  \fv{oef:h1-t4}{vb:h1-grondformule}{oef:h1-grondf}%
  \fv{oef:h1-t5}{vb:h1-verwant}{oef:h1-herleid}%
  \fv{oef:h1-t6}{vb:h1-waarden}{oef:h1-punt}%
  \fv{oef:h1-t7}{vb:h1-verwant}{oef:h1-verwant}%
  \fzv{oef:h1-t8}{oef:h1-o6}}

\ifzelfstudie
\section*{Maak je eigen samenvatting}
Sluit dit boek en schrijf in je schrift: de definitie van een radiaal, het omrekenen, booglengte en sectoroppervlakte, de definities van $\sin$, $\cos$ en $\tan$ op de eenheidscirkel, de grondformule, de tabel met bijzondere waarden en de vier soorten verwante hoeken, telkens met een figuur. Vergelijk daarna met de kaders.
\fi
